\section{Introduction}
\label{sec:intro}

Recovering an image from the magnitude of its Fourier transform is a classical, ill-posed inverse problem~\cite{fienup1982phase,shechtman2015phase}.
Diffusion models now supply the image prior for many solvers of this task~\cite{chung2023dps,zhang2025daps,li2024difffpr,alkhouri2025sitcom,kim2026fastdips}, and their reported quality has risen quickly.
The same papers also note that individual runs fail often: the reconstruction may be rotated, color-scrambled, or unrelated to the image~\cite{zhang2025daps,kocabas2026pcdaps}.
The field's response has been to run each solver several times and report the best run.
In the reference implementations of DAPS, SITCOM and FAST-DIPS, ``best'' is decided against the ground truth and for every metric separately, so the reported PSNR, SSIM and LPIPS may come from three different runs~\cite{zhang2025daps,alkhouri2025sitcom,kim2026fastdips}\claim{C-P1}.

Two recent works question this practice.
Kocabas~\etal\cite{kocabas2026pcdaps} show that ground-truth best-of-$k$ masks frequent orientation failures and prune particles by prior consistency instead; Makris~\etal\cite{makris2026tgd} select runs by their measurement error and remark that the per-metric oracle is not available at test time.
Measurement error is the natural ground-truth-free criterion, but it is not obviously adequate: phase retrieval is ill-posed, so a small residual does not certify a correct image.
Solvers already add prior or network-consistency terms \emph{inside} the reconstruction for this reason~\cite{alkhouri2025sitcom}; whether the \emph{selection} among finished runs needs them too, and how much of the published quality is available without the ground truth, has not been measured.

We find that measurement-based selection fails in exactly three ways, with different causes and remedies.
(i)~It is provably blind to an exact eight-fold symmetry of the operator, so rotated runs pass; a cheap prior or semantic test fixes them and is the Bayes rule within the orbit (\cref{sec:symmetry}).
(ii)~It cannot reject runs that fit the measurement within the noise but show a different image; how often these occur depends strongly on the solver, and their share of selection errors grows with the number of runs (\cref{sec:selection,sec:audit}).
(iii)~It cannot help when no run is close to the image, which is the dominant failure at the run counts used in practice\claim{C-A4}.
On FFHQ, residual selection followed by an orientation test recovers $93$--$97\%$ of the gain of an orientation-aware ground-truth oracle\claim{C-A3}; what remains is mostly~(iii).

\paragraph{The symmetry behind (i).}
The oversampled Fourier operator used by these solvers is exactly invariant under eight transformations of the image: independent $180^\circ$ rotations of each color channel (\cref{fig:teaser}a).
This ambiguity is classical~\cite{hayes1982reconstruction,fienup1986stagnation} and was observed for diffusion solvers by Zhang~\etal\cite[App.~F]{zhang2025daps}.
We turn it into precise statements.
Any selection rule or particle weight computed from the measurements is blind to the orientation of a run (\cref{sec:symmetry}).
And if the denoiser were equivariant to these rotations, every guided sampler we study would choose among the eight orientations uniformly at random; the strong preference for upright images that we observe must therefore come from the network alone.
We bound this preference by the network's equivariance defect, which we measure for an FFHQ and an ImageNet denoiser.

\paragraph{A model of best-of-\texorpdfstring{$K$}{K}.}
Selecting the run with the smallest measurement residual is resampling with an imperfect verifier~\cite{stroebl2026limits}: runs of type~(ii) are false positives, and rotated runs are false positives that no measurement-based verifier can reject.
We derive closed-form curves for the success of ground-truth-free and ground-truth selection as a function of $K$, exact estimators of these curves from a finite set of runs, and the fact that the gap between them never shrinks with $K$ (\cref{sec:selection}).
We fit the model on four runs per image and freeze its predictions before collecting up to $64$ runs per image.

\paragraph{An audit.}
We then decompose the reported quality of three published solvers (DAPS, SITCOM, FAST-DIPS; DPS, PC-DAPS and DiffFPR in the supplement \pending{if run}) on FFHQ and ImageNet into the solver itself, ground-truth-free selection, the ground-truth oracle, and the per-metric oracle (\cref{sec:audit}).

\paragraph{Contributions.}
\begin{itemize}[leftmargin=*,itemsep=1pt,topsep=2pt]
\item \textbf{Symmetry theory.} The orbit of the standard operator; the blindness of measurement-based selection and of likelihood-weighted particle methods to it; and a total-variation bound tying a sampler's orientation preference to the denoiser's equivariance defect, which we measure (\cref{sec:symmetry}).
\item \textbf{A selection model} for best-of-$K$ in inverse problems, extending the imperfect-verifier model of~\cite{stroebl2026limits} with orbit classes; closed-form $K$-curves, exact estimators, and predictions frozen before the data and tested up to $K{=}64$ (\cref{sec:selection}).
\item \textbf{An audit} of three published solvers on two datasets and three metrics that separates each layer of the protocol and the three failure types (\cref{sec:audit}).
\item \textbf{Reporting recommendations} (\cref{sec:discussion}).
\end{itemize}
